\documentclass[11pt]{article}
\usepackage[letterpaper,margin=0.75in]{geometry}
\usepackage[T1]{fontenc}
\usepackage{lmodern}
\usepackage{amsmath,amssymb,amsthm,microtype,booktabs,xcolor}
\usepackage[colorlinks=true,urlcolor=blue!45!black,linkcolor=blue!45!black,citecolor=blue!45!black]{hyperref}
\newtheorem{theorem}{Theorem}
\newtheorem{lemma}[theorem]{Lemma}
\newcommand{\G}{\operatorname{G}}
\newcommand{\xor}{\mathbin{\mathrm{xor}}}
\newcommand{\floor}[1]{\left\lfloor #1\right\rfloor}
\setlength{\parindent}{0pt}
\setlength{\parskip}{0.35em}
\hypersetup{pdftitle={The growth rate and alternating ratios of OEIS A265386},pdfauthor={Blake Stone}}
\begin{document}
\begin{center}
{\Large\bfseries The growth rate and alternating ratios\\[0.15em]of OEIS A265386}\\[0.65em]
Blake Stone\\[0.15em]
{\small Proof developed with OpenAI Astra 6.2\quad$\vert$\quad November 26, 2026}
\end{center}
\begin{quote}\small\textbf{Abstract.}
We prove Sykora's conjectured mean growth rate $\sqrt2$ for the recurrence obtained by applying binary-reflected Gray code to both inputs and the output of a Fibonacci step. An explicit set of 48 binary-prefix states is closed under all possible carries. This finite certificate, together with induction, bounds every later term. It also proves that adjacent ratios alternate between values greater than $3$ and values less than $2/3$. Both results are checked in Lean.
\end{quote}

\section*{1. The recurrence and the prefix transition}
For a nonnegative integer $m$, write
\[
 \G(m)=m\xor\floor{m/2},\qquad F(u,v)=\G\bigl(\G(u)+\G(v)\bigr).
\]
Define $a(1)=a(2)=1$ and
\[
 a(n)=F\bigl(a(n-1),a(n-2)\bigr)\qquad(n\ge3).
\]
This is OEIS A265386, contributed by Stanislav Sykora in 2015~\cite{oeis}.

\begin{theorem}\label{thm:main}
The mean growth rate is
\[
 \boxed{\displaystyle\lim_{n\to\infty}\left(\frac{a(2n)}{a(n)}\right)^{1/n}=\sqrt2.}
\]
Moreover, for every integer $k\ge8$,
\[
 \frac{a(2k+1)}{a(2k)}>3,
 \qquad
 \frac{a(2k+2)}{a(2k+1)}<\frac23.
\]
In particular, the adjacent-term ratios do not converge.
\end{theorem}

\textbf{Prefix identity.} For all nonnegative integers $m,r$,
\begin{equation}\label{eq:gray}
 \floor{\frac{\G(m)}{2^r}}=\G\left(\floor{\frac m{2^r}}\right).
\end{equation}
Indeed, if $b_i$ is bit $i$ of $m$, bit $i$ of $\G(m)$ is $b_i\xor b_{i+1}$. Deleting the lowest $r$ bits before or after the transform therefore gives the same result.

\textbf{Two-step transition.} Suppose $P=2^r$, $x=\lfloor u/P\rfloor$, and $y=\lfloor v/P\rfloor$. Quotients of sums obey
\[
 \floor{\frac{s+t}{P}}=\floor{\frac sP}+\floor{\frac tP}+c,
 \qquad c\in\{0,1\}.
\]
Apply this fact and \eqref{eq:gray} to the next two recurrence steps. For some $c,d\in\{0,1\}$, define
\begin{equation}\label{eq:transition}
 z=\G\bigl(\G(x)+\G(y)+c\bigr),\quad
 w=\G\bigl(\G(y)+\G(z)+d\bigr),\quad
 T_{c,d}(x,y)=(\lfloor z/2\rfloor,\lfloor w/2\rfloor).
\end{equation}
Then the leading quotients of $F(u,v)$ and $F(v,F(u,v))$ at scale $2P$ are exactly $T_{c,d}(x,y)$. The actual carries need not be independent; allowing all four combinations includes every possibility.

\newpage
\section*{2. A finite certificate for all indices}
Let $S$ be the following set of 48 ordered pairs. Each row lists one first coordinate and all corresponding second coordinates.
\begin{center}
\renewcommand{\arraystretch}{1.12}
\begin{tabular}{rl}
\toprule
$x$ & All $y$ with $(x,y)\in S$\\
\midrule
136 & 504, 505, 506, 507\\
137 & 506, 507, 510\\
138 & 497, 499, 506, 508, 510\\
139 & 509, 510, 511\\
142 & 489, 491, 500\\
148 & 496, 497\\
152 & 502\\
154 & 497, 498, 499, 502, 503\\
155 & 501, 502, 503\\
156 & 506, 507, 510, 511\\
157 & 506, 508, 509, 510, 511\\
158 & 500, 506, 507, 510, 511\\
159 & 504, 505, 506, 507, 510\\
\bottomrule
\end{tabular}
\end{center}
Direct integer calculation verifies the finite closure property
\begin{equation}\label{eq:closed}
 (x,y)\in S,\quad c,d\in\{0,1\}
 \quad\Longrightarrow\quad T_{c,d}(x,y)\in S.
\end{equation}
There are exactly $48\cdot4=192$ cases. The accompanying Lean source spells out $S$ and proves \eqref{eq:closed} by kernel evaluation of this finite proposition. All states also satisfy
\begin{equation}\label{eq:coordinate}
 136\le x\le159,\qquad489\le y\le511.
\end{equation}

\begin{lemma}[Invariant]\label{lem:invariant}
For every integer $k\ge8$,
\begin{equation}\label{eq:invariant}
 \left(\floor{\frac{a(2k)}{2^{k-8}}},\floor{\frac{a(2k+1)}{2^{k-8}}}\right)\in S.
\end{equation}
\end{lemma}
\begin{proof}
Direct evaluation of the original recurrence gives $a(16)=158$ and $a(17)=510$. Since $(158,510)\in S$, the statement holds for $k=8$. If it holds at $k$, the two-step transition in \eqref{eq:transition} and the closure property \eqref{eq:closed} give the next pair at doubled scale, which is precisely \eqref{eq:invariant} at $k+1$. Induction completes the proof.
\end{proof}

Set $P=2^{k-8}$. The invariant and \eqref{eq:coordinate} imply the useful bounds
\begin{align}
 136P&\le a(2k)<160P,\label{eq:even}\\
 489P&\le a(2k+1)<512P.\label{eq:odd}
\end{align}
They prove positivity of every term from index $16$ onward. They also imply the weaker bounds
\[
 2^{k-1}\le a(2k)<2^k,
 \qquad 2^k\le a(2k+1)<2^{k+1}.
\]
After squaring and separating the two parities, we obtain the uniform estimate
\begin{equation}\label{eq:squared}
 \boxed{2^n\le4a(n)^2,\qquad a(n)^2<2\cdot2^n\qquad(n\ge16).}
\end{equation}

\newpage
\section*{3. Growth and persistent oscillation}
\begin{proof}[Proof of Theorem~\ref{thm:main}]
Apply \eqref{eq:squared} at $n$ and $2n$. Positivity gives
\[
 \frac18\,2^n
 <\left(\frac{a(2n)}{a(n)}\right)^2
 <8\,2^n\qquad(n\ge16).
\]
Taking positive real $2n$th roots yields
\[
 \sqrt2\,8^{-1/(2n)}
 <\left(\frac{a(2n)}{a(n)}\right)^{1/n}
 <\sqrt2\,8^{1/(2n)}.
\]
For every fixed $c>0$, $c^{1/n}\to1$. The squeeze theorem proves the conjectured limit $\sqrt2$.

For the adjacent ratios, let $k\ge8$ and $P=2^{k-8}$. Equations \eqref{eq:even}--\eqref{eq:odd} give
\[
 3a(2k)<480P<489P\le a(2k+1).
\]
The upper bound at $k+1$ is $a(2k+2)<320P$, so
\[
 3a(2k+2)<960P<978P\le2a(2k+1).
\]
Dividing by the positive denominators proves both claimed ratio inequalities. If adjacent ratios had a common limit, the even-start subsequence would force it to be at least $3$, while the odd-start subsequence would force it to be at most $2/3$, a contradiction.
\end{proof}

\section*{4. Formal verification}
The package pins Lean 4.19.0 and a matching Mathlib revision. The module \texttt{Proof.lean} imports only \texttt{Std}: it defines the exact recurrence and 48 states, checks all 192 transitions by ordinary \texttt{decide}, and proves the invariant, parity bounds, positivity, and adjacent-ratio inequalities. The module \texttt{Limit.lean} proves \eqref{eq:squared} and the unconditional real limit theorem \texttt{A265386.growth\_rate\_limit}. Its \texttt{meanGrowth} definition is exactly $(a(2n)/a(n))^{1/n}$, using real division and exponentiation.

Both packaged modules passed a fresh build and an independent rebuild. Their audited theorems use only the standard logical dependencies \texttt{propext}, \texttt{Classical.choice}, and \texttt{Quot.sound}. There are no unfinished proofs, added axioms, or \texttt{native\_decide} calls. The finite closure certificate is connected to every index by induction. Reproduction commands and build output are supplied in \texttt{README.md}, \texttt{verification.txt}, and the independent audit record.

\section*{Source status}
The live OEIS entry inspected on November 26, 2026 still described the mean growth limit as a conjecture; its internal revision was \#16, January 16, 2016. Searches by identifier and recurrence description, including available formal-proof collections, found no earlier proof of this assertion. This records the source status checked during the work and does not establish exhaustive publication priority.

\begin{thebibliography}{9}
\bibitem{oeis} S. Sykora, \emph{A265386}, The On-Line Encyclopedia of Integer Sequences, December 7, 2015. \url{https://oeis.org/A265386}. Accessed November 26, 2026.
\end{thebibliography}
\end{document}
